% SPDX-FileCopyrightText: Copyright (c) 2023-2025 Yegor Bugayenko
% SPDX-License-Identifier: MIT

\documentclass[nobrand,anonymous,nodate,nosecurity]{huawei}
\usepackage{enumerate}
\usepackage{bibentry}
\usepackage{multicol}
\usepackage{href-ul}
\usepackage[noframes]{ffcode}
\usepackage{soul}
\usepackage{iexec}
\usepackage{booktabs}
\usepackage{svg}
\bibliography{../bibliography/main.bib}

\begin{document}

\iexec[quiet,log]{wget --no-check-certificate -N --quiet https://www.yegor256.com/images/face-1024x1024.jpg}
\iexec[quiet,log]{wget --no-check-certificate -N --quiet https://www.objectionary.com/cactus.svg}

\includesvg[height=5em]{cactus.svg}

\section*{Pain of OOP}

This series of lectures by \href{https://www.yegor256.com}{Yegor Bugayenko} was first presented to students at \href{https://innopolis.university/en/}{Innopolis University} in 2023 and was \href{https://www.youtube.com/playlist?list=PLaIsQH4uc08ytf8POIIAkkR4ZsRq8DFiV}{video-recorded}.
The complete set of slide decks is hosted and maintained in the \href{https://github.com/yegor256/painofoop}{yegor256/painofoop} GitHub repository.

\textbf{What is the course about?}\\
This course offers a critical review of the current state of object-oriented programming, with a particular focus on the Java, C++, Ruby, and JavaScript ecosystems.
It examines various programming idioms—often referred to as ``best practices''---and critiques their negative impact on code quality.
These include static methods, null references, getters and setters, ORM and DTO patterns, annotations, traits and mixins, inheritance, and others.
In contrast, the course proposes alternative approaches that promote cleaner object-oriented programming practices.

\textbf{What is the goal?}\\
The primary objective of the course is to help students understand the difference between ``object thinking''---which originally motivated the emergence of object-oriented programming---and modern practices that often severely degrade code quality.

\floatstyle{plain}
\restylefloat{figure}
\begin{wrapfigure}{r}{5em}%
\raggedleft%
\includegraphics[width=\linewidth]{face-1024x1024.jpg}%
\end{wrapfigure}
\textbf{Who is the teacher?}\\
Yegor has been developing software for over 30 years, working both as a hands-on programmer (see his GitHub account: \href{https://github.com/yegor256}{@yegor256}) and as a manager of other programmers.
He is currently a software developer at Huawei (Moscow, Russia).
His primary research focus is on software quality issues.
Some of the lectures he has recently presented at software conferences can be found on \href{https://www.youtube.com/channel/UCr9qCdqXLm2SU0BIs6d_68Q}{his YouTube channel}.
Yegor has also published \href{https://www.yegor256.com/books.html}{several books} and maintains a \href{https://www.yegor256.com/contents.html}{blog} about software engineering and object-oriented programming.
He has previously taught courses at Innopolis University (Kazan, Russia) and HSE University (Moscow, Russia), including \href{https://github.com/yegor256/ppa}{Practical Program Analysis (2023)} and \href{https://github.com/yegor256/sqm}{Software Quality Metrics (2024)} (all videos are available).

\textbf{Why this course?}\\
The maintainability of object-oriented software that most of us programmers write today falls far short of our expectations.
One of the main reasons for this is a fundamental misunderstanding of what \emph{objects} truly are.
This course may help clear things up.

% \textbf{What's the methodology?}\\
% Each lecture is a critical review of an existing ``best practice,'' such as static methods or getters, with the intent to highlight its negative impact on software quality.

\newpage
\section*{Course Structure}

Prerequisites to the course (it is expected that a student knows this):

\begin{itemize}
\item How to write code
\item How to design software
\item How to use Git and GitHub
\end{itemize}

After completing the course, a student will \ul{hopefully} understand:

\begin{itemize}
\item What is the difference between objects and data?
\item Why are static methods bad?
\item What is immutability and why is it good?
\item How should one design a constructor?
\item How can exceptions be handled correctly?
\item What is data hiding for?
\item What's wrong with Printers, Writers, Scanners, and Readers?
\item Why are NULL references considered a billion-dollar mistake?
\item What's wrong with mixins and traits?
\item How can getters and setters be avoided?
\item How is declarative programming better than imperative programming?
\item Why is composition better than inheritance?
\item Where should data be stored if DTOs are considered bad practice?
\item How can the ORM design pattern be avoided?
\item Why are long variable names considered bad design?
\item Why are type casting and type checking against OOP principles?
\item What is cohesion and why does it matter?
\item How can SOLID and SRP principles be applied?
\item What is Inversion of Control for, and why are DI Containers considered harmful?
\item Why is MVC considered a bad design idea?
\end{itemize}

\newpage
\section*{Lectures}

The following topics are discussed:

\newlist{lectures}{enumerate}{10}
\setlist[lectures]{label*=\arabic*.}
\begin{lectures}
\item Algorithms
\item Static Methods
\item Getters
\item Setters and Mutability
\item ``-ER'' Suffix
\item NULL References
\item Type Casting and Reflection
\item Inheritance
\end{lectures}

\newpage
\section*{Grading}

In order to pass the course, students must attend lectures and labs.
They must also contribute to one of the following GitHub repositories, which are written in more or less ``pure'' object-oriented style:

\begin{itemize}
    \item \href{https://github.com/objectionary/eo}{\ff{objectionary/eo}} (Java + XSLT)
    \item \href{https://github.com/yegor256/cactoos}{\ff{yegor256/cactoos}} (Java)
    \item \href{https://github.com/yegor256/takes}{\ff{yegor256/takes}} (Java)
    \item \href{https://github.com/yegor256/factbase}{\ff{yegor256/factbase}} (Ruby)
    \item \href{https://github.com/objectionary/eoc}{\ff{objectionary/eoc}} (JavaScript)
    \item \href{https://github.com/objectionary/eo2js}{\ff{objectionary/eo2js}} (JavaScript)
    \item \href{https://github.com/yegor256/cam}{\ff{yegor256/cam}} (Bash + Python)
\end{itemize}

There is no oral or written exam at the end of the course.
Instead, each student earns points for the following results:

\renewcommand{\arraystretch}{1}
\begin{tabular}{lrr}
\toprule
Result & Points & Limit \\
\midrule
Attended a lecture & +1 & 6 \\
Attended a lab & +1 & 8 \\
Submitted a ticket (that was accepted) & +2 & 8 \\
Submitted a pull request (that was merged) & +4 & 32 \\
\bottomrule
\end{tabular}

Final grades are based on accumulated points: 25+ earns an ``A,'' 17+ a ``B,'' and 9+ a ``C.''

If eight or fewer points are accumulated during the course, the student may still earn them before the retake.

An online lecture is counted as ``attended'' only if a student was personally present in Zoom for more than 75\% of the lecture's time.
Watching the lecture from the computer of a friend \ul{doesn't} count.

\newpage
\renewcommand\refname{Learning Materials}
\bibentry{beck1997smalltalk}
\bibentry{booch1994object}
\bibentry{eo1}
\bibentry{eo2}
\bibentry{martin2008cleancode}
\bibentry{martin2017clean}
\bibentry{mcconnell2004codecomplete}
\bibentry{pawson2002naked}
\bibentry{west2004object}
\printbibliography

You may also want to explore Yegor's blog posts, \href{https://www.yegor256.com/tag/oop}{on his blog}, his video lectures \href{https://www.youtube.com/playlist?list=PLaIsQH4uc08yw2CsNv5OV30GfKE6XVGii}{on YouTube}, and Object Thinking meetup presentations \href{https://www.youtube.com/watch?v=yT6oO28wEik&list=PLaIsQH4uc08yetzX86w1pPck1QtGEy_ik}{on YouTube}.

\end{document}
